\section{Scope and Limitations}
\label{sec:limitations}
The controlled design isolates representation and capacity within a fixed
short-window setting, but its scope is intentionally bounded. The main analysis
uses one synthetic corpus, a 128-sample observation, one primary backbone, a
predefined capacity range, and a finite set of seeds and data splits. A small
secondary-backbone control is used for directional checking rather than as a
second full benchmark. The study reports parameter and operation-count costs;
it does not make a physical-device latency or energy claim. The split audit
reports the source metadata available for the corpus.

The present results use nominal width as the capacity axis. Although the cost
ledger stores Core FLOPs and can define admissible matching pairs, no
Core-FLOP-matched accuracy comparison is reported here. Likewise, the
pre-specified H4 SNR-region analysis and H5 secondary-backbone control are not
included in the current results. They should not be inferred from the nominal
width contrasts.

All widths share the same learning-rate, weight-decay, and regularization
settings. The reported capacity effects therefore describe the common protocol
and do not establish that each width was individually hyperparameter-optimal.
The nested feature ladder also changes both information content and the number
of input channels, so it does not identify which individual derived channel
causes an observed effect. Finally, the study cannot infer source-record
independence when the dataset metadata do not expose waveform provenance.
The analysis plan also specified a split-aware mixed model, but it was not fit
because only two split levels were available. The manuscript therefore uses
split-stratified descriptive checks and paired split--seed summaries instead.
